\documentclass[11pt]{scrartcl}

\usepackage[stixtwo]{fontsetup}
\usepackage{amsmath}

\begin{document}
\begin{enumerate}
	\item Például legyen a sorozat
	      \[
		      1;\ 1, 2;\ 1, 2, 3;\ 1, 2, 3, 4;\ \ldots
		      \text,
	      \]
	      vagyis az $m$-edik blokk az $1, 2, \ldots, m$ számokból áll.

	      Minden $k$ pozitív egész szám minden $m \ge k$ indexű blokkban
	      szerepel, tehát végtelen sokszor fordul elő a sorozatban, így
	      torlódási pont.

	      Más valós szám nem lehet torlódási pont, hiszen a sorozat elemei
	      pozitív egészek: ha $c$ nem pozitív egész, akkor $c$-nek van olyan
	      környezete, amely egyetlen sorozatelemet sem tartalmaz.

	\item \textbf{Bizonyítás}

	      A $+\infty$ akkor torlódási pont, ha minden $K$ valós számhoz
	      végtelen sok olyan $n$ index létezik, amelyre $a_n > K$.

	      Legyen $K$ tetszőleges, és válasszunk egy $k > K + 1$ pozitív egész
	      számot. Mivel $k$ torlódási pontja a sorozatnak, ezért végtelen sok
	      $n$-re
	      \[
		      |a_n - k| < 1
		      \quad \Longrightarrow \quad
		      a_n > k - 1 > K
		      \text.
	      \]

	      Következésképpen: a sorozatnak bármely $K$ felett végtelen sok eleme
	      van, vagyis a $+\infty$ is torlódási pontja.
\end{enumerate}
\end{document}
